% ECONOMICS_PROBLEM: {"id": "C73-445", "title": "Convergence of model-free learning in dynamic population games", "jel_code": "C73", "source_id": "OP-0152"}
% !TeX program = xelatex
\documentclass[11pt,a4paper]{article}
\usepackage[margin=23mm]{geometry}
\usepackage{fontspec}
\setmainfont{Latin Modern Roman}
\usepackage{amsmath,amssymb,mathtools}
\usepackage{xcolor,enumitem,booktabs,longtable,array,xurl,hyperref}
\definecolor{muted}{HTML}{53636C}
\hypersetup{colorlinks=true,urlcolor=blue,linkcolor=blue}
\setlength{\parindent}{0pt}
\setlength{\parskip}{5pt}
\newcommand{\R}{\mathbb R}
\newcommand{\E}{\mathbb E}
\newcommand{\N}{\mathbb N}
\newcommand{\ind}{\mathbf 1}
\newcommand{\Var}{\operatorname{Var}}
\newcommand{\Cov}{\operatorname{Cov}}
\newcommand{\supp}{\operatorname{supp}}
\DeclareMathOperator*{\argmin}{arg\,min}
\DeclareMathOperator*{\argmax}{arg\,max}
\newcommand{\entrymeta}[3]{{\sffamily\small\color{muted}\textbf{\texttt{#1}}\quad|\quad #2\par #3\par}}
\newcommand{\Problemref}[1]{\texttt{#1}}
\newcommand{\JELref}[2]{\texttt{#2} (JEL \texttt{#1})}
\begin{document}
\section*{Convergence of model-free learning in dynamic population games}
\textbf{Problem ID:} \texttt{C73-445}\quad \textbf{JEL:} \texttt{C73}\par
% BEGIN ECONOMICS STATEMENT
\label{problem:OP-0152}
% GAME_THEORY_PROBLEM: {"local_id": "GT-022", "original_number": 22, "kind": "R", "entry_kind": "statement_only", "published": false, "review_required": true, "category": "game_theory"}
\label{game:GT-022}
{\sffamily\small\color{muted}
{\sffamily\small\color{muted}\textbf{GT-022} | Learning in games\par R | Implementation-dependent convergence agenda\par}
\par}
\subsection*{Definitions and mathematical statement}
Let $X$ be a finite state space and $A_x$ a nonempty finite action set at state $x$. A population configuration is $z=(\mu,\pi)$, with $\mu\in\Delta(X)$ and $\pi(\cdot\mid x)\in\Delta(A_x)$. For each $z$, specify bounded rewards $r_z(x,a)$, transition probabilities $p_z(y\mid x,a)$, and a discount factor $\alpha\in[0,1)$. Let $V_z^*$ be the unique solution of
\[
 V_z^*(x)=\max_{a\in A_x}\left(r_z(x,a)+\alpha\sum_{y\in X}p_z(y\mid x,a)V_z^*(y)\right).
\]
The stationary Nash configuration set $\mathcal E$ comprises configurations satisfying
\begin{align*}
 \mu(y)&=\sum_{x\in X}\sum_{a\in A_x}\mu(x)\pi(a\mid x)p_z(y\mid x,a),\qquad y\in X,\\
 \pi(a\mid x)>0&\ \Longrightarrow\ 
 r_z(x,a)+\alpha\sum_y p_z(y\mid x,a)V_z^*(y)=V_z^*(x).
\end{align*}
This uses the optimal value for the environment frozen at $z$, rather than assuming that a value estimated under an arbitrary current policy is already optimal.

For a completely specified stochastic implementation $\mathcal A_\theta$ producing $z_k$, a precise convergence proposition is
\[
 \mathcal E\ne\varnothing\quad\Longrightarrow\quad
 \Pr\!\left(\lim_{k\to\infty}\operatorname{dist}(z_k,\mathcal E)=0\right)=1,
\]
where distance is Euclidean distance on the finite product of simplices. The research target is to establish such a proposition, or exhibit failure, for the source's noncontractive dynamic-population models and its FP-DQN procedure.
\subsection*{Short English statement}
Does the proposed model-free learning process approach the set of stationary population equilibria in the intended models?
\subsection*{Variants and scope boundaries}
The implementation parameter $\theta$ must fix the network class and approximation error, optimizer, learning rates, exploration, replay-buffer rule, sample/population sizes, and their asymptotic schedules. With persistent finite-sample noise, an error neighborhood or convergence in distribution may be the correct target. Convergence to one selected equilibrium is stronger than distance-to-set convergence when equilibria are not unique.
\subsection*{Source relationship and status}
[S1] supplies Algorithm 1 and numerical evidence, not a general convergence theorem without contraction or uniqueness. Sections 5 and 7 discuss these limitations. The catalog sentence alone does not fix $\theta$ sufficiently to define a single stochastic process; the displayed proposition states exactly what remains to be specified before proving or disproving it.
\subsection*{References}
\begingroup\footnotesize\raggedright
\textbf{[S1]} Matteo Cederle, Saverio Bolognani and Gian Antonio Susto, \emph{Towards Model-Free Learning in Dynamic Population Games: An Application to Karma Economies}, arXiv:2605.11042v1 (2026). Locators: Section 3 (stationary Nash configuration), Section 5 and Algorithm 1 (FP-DQN and limitations), Section 7 (concluding discussion). \url{https://arxiv.org/html/2605.11042v1}
\par\endgroup

\subsection*{Common conventions: Source mathematical conventions}

\textbf{Finite games.} For a finite set $A$, $\Delta(A)$ is its probability
simplex. Players choose independent mixed strategies in Nash equilibrium;
correlation is allowed only when explicitly part of the solution concept.
In a game with payoff functions $u_i$ and strategy profile $x$, an additive
$\varepsilon$ Nash equilibrium satisfies
\[
 u_i(x)\geq u_i(x_i',x_{-i})-\varepsilon
 \quad\text{for every player $i$ and every admissible deviation $x_i'$}.
\]
For numerical approximation, payoffs are normalized as locally specified.
An $\varepsilon$ payoff residual is distinct from distance at most
$\varepsilon$ to the set of exact equilibria. Point convergence, convergence
of empirical play, and convergence in distance to an equilibrium set are
also distinct.

\textbf{Computational representations.} Unless stated otherwise, a complexity
question takes rational data with binary encoding; polynomial time is measured
in total bit length. A fixed-accuracy polynomial algorithm is not necessarily
polynomial in $1/\varepsilon$, and neither is necessarily polynomial in
$\log(1/\varepsilon)$. Succinct games and value, demand, or sampling oracles
must declare their interfaces. Exact real solutions require an explicit
representation; an unspecified list of arbitrary real numbers is not a
finite computational output. A claimed algorithmic barrier is meaningful
only for its specified model and reductions.

\textbf{Dynamic games.} Histories, information, observation, and allowable
randomization are part of the model. A stationary strategy depends only on
the current state; a positional strategy in a deterministic graph game
chooses one action at each controlled vertex. Nash equilibrium at an initial
state and equilibrium after every history are different requirements.
An expectation of a pathwise $\liminf$ payoff, a $\liminf$ of expected averages,
a limiting expected average, and a uniform finite-horizon equilibrium are
different criteria. None is silently substituted for another.

\textbf{Allocation and mechanisms.} A complete allocation assigns every item
exactly once unless otherwise stated. Zero-valued goods, negative values,
capacity constraints, feasibility, lotteries, and copies of goods affect
fairness definitions and are specified locally. Dominant-strategy incentive
compatibility, Bayesian incentive compatibility, universal truthfulness,
and truthfulness in expectation impose different inequalities. Whenever
an objective ratio has zero denominator, its convention is stated or those
instances are excluded.

\textbf{Infinite-dimensional models.} Expectations, integrals, stochastic
equations, and optimization objectives require their local regularity,
integrability, filtrations, and admissible domains. A backward--forward
mean-field system is a boundary-value problem; it is not an initial-value
semiflow without an additional construction. Long-time, large-population,
and vanishing-noise limits require an order of limits and a convergence
topology. If a minimum need not be attained, the objective is its infimum
or supremum and the approximation requirement is stated separately.
% END ECONOMICS STATEMENT
\end{document}
